\begin{lemma}[A positive-divider minimum hierarchy counterexample cannot have one tail addition]
Fix a finite loopless multigraph with its individual edge identities, a
palette of $k$ colors, and a distinguished edge $e_0$.  Let
$(T,\varphi_n)$ be a positive-rung extended Tashkinov tree constructed
from an arbitrary certified input coloring, let $\mathcal B$ be a
hierarchy base under $\varphi_n$, and let $C$ be a hierarchy
counterexample over that base.  Write $H=C.\mathit{hierarchy}$, let $q$
be its divider count, and put
\[
 A=\noderef{HierarchyStarredStage}(H,q),\qquad
 \ell=|T|-|A|,
\]
and let
\[
 T(y_{\ell-1})=
 \noderef{TreeSequencePrefix}
   \bigl(T,\min\{|A|+(\ell-1),|T|\}\bigr).
\]
Assume that $\varphi_n$ is proper and leaves exactly $e_0$ uncolored.

Assume, uniformly for every positive-rung maximum ETT on the same graph,
palette, and distinguished edge, with arbitrary certified input and
generating colorings and with rung number smaller than that of $T$, all
six lower-rung assertions of the technical theorem: elementarity;
terminal interchangeability; at every PE slot, unconditional stability
before the support, stability on the preceding ladder tree when the
first-appearance index is greater than one, and disjointness of stored
pairs at distinct PE slots with distinct supports and equal
first-appearance index; the terminal PE component assertion; the
defective-endpoint order assertion; and the mod-coloring assertion.
Assume that every same-rung maximum ETT with a good hierarchy having
fewer than $q$ dividers is elementary.

Assume that $C$ is lexicographically minimal among hierarchy
counterexamples whose ETTs have the same rung number as $T$ and whose
hierarchies have exactly $q$ dividers, allowing the certified input
coloring, generating coloring, ETT, base, and hierarchy to vary.  Thus
no such $C'$ has
\[
 \noderef{HierarchyCounterexampleLess}
   (\noderef{HierarchyCounterexampleMeasure}(C'),
    \noderef{HierarchyCounterexampleMeasure}(C)).
\]
Assume also the large-branch estimate
\[
 2n+11\leq |\overline\varphi_n(T_n)|,
\]
where $n$ is the rung number and $T_n$ is the terminal ladder tree.

If
\[
 \ell=|T|-|A|=1,\qquad
 V(T(y_{\ell-1}))\text{ is elementary under }\varphi_n,
 \qquad 0<q,
\]
then a contradiction follows.
\end{lemma}
\begin{proof}
All tree inclusions and equalities below concern complete ordered
edge-valued tree-sequences.  Since $0<q$, the positive-index clause of
\noderef{HierarchyStarredStage} gives $A=H.\mathit{stage}(q)$.
The nested, strict, and terminal-stage fields of
\noderef{TashkinovHierarchy}, together with $\ell=1$, give
$A\prec T$ and $|T|=|A|+1$.  The entrywise clauses of
\noderef{TreeSequenceInitialSegment} therefore display the exact final
addition as
\[
                         T=A\cup\{e_1,y_1\}.                  \tag{1}
\]
The edge $e_1$ is the actual edge entry at index $|A|$, not an endpoint
pair.  Since $\ell-1=0$, the assumed
\noderef{TreeSequencePrefix} has length $|A|$ and is entrywise equal to
$A$.  Hence
\[
              V(A)\text{ is elementary under }\varphi_n.     \tag{2}
\]

Expand the `notElementary' field of
\noderef{HierarchyCounterexample} using
\noderef{ColoredVertexSetElementary}.  It gives distinct
$u,v\in V(T)$ and a common missing color $\alpha$.  By
\noderef{TreeSequenceVertexSet} and (1),
$V(T)=V(A)\cup\{y_1\}$.  Equation (2) excludes two such vertices in
$A$, while distinctness excludes two copies of $y_1$.  After exchanging
$u,v$ if needed and expanding \noderef{TreeMissingColors} and
\noderef{MissingColors},
\[
 \alpha\in\overline\varphi_n(A)
       \cap\overline\varphi_n(y_1).                           \tag{3}
\]

Write $n$ for the rung number, $T_n$ for the terminal ladder tree, and
$T_{n+1}$ for the last output tree of the certified series.  The
`certified' field of \noderef{TashkinovSeries}, evaluated at the actual
last index $\operatorname{Fin.last}(n)$, supplies a
\noderef{TashkinovSeriesStep} certificate at that index.  The
`roundsPositive' field of \noderef{ExtendedTashkinovTree} says $n\geq1$,
so this is a positive index and the certificate cannot be the initial
constructor.  It is therefore an actual SE, PE, or RE constructor.  The
other fields of \noderef{ExtendedTashkinovTree} give the exact ordered
relations and coloring identity
\[
                 T_n\prec T\preceq T_{n+1},
        \qquad \varphi_n=\varphi^{\rm series}_n.              \tag{4a}
\]
In the PE constructor the certified final
\noderef{TashkinovClosure} starts at $T_n$.  In either the SE or RE
constructor it records an individual edge identity $f_n$, the singleton
connecting-set equality $F_n=\{f_n\}$, the exact seed $T_n+f_n$ in the
sense of \noderef{TreeSequenceAdjoinEdge}, and a certified closure from
that seed to $T_{n+1}$.  By
\noderef{HierarchyNonparallelFirstEdge}, this same $f_n$ is the first
edge of the first positive hierarchy stage.  Since $q>0$, that stage is
an ordered prefix of $A$.  The exact seed and that stage are also ordered
prefixes of the same closure output; the edge and vertex entry clauses of
\noderef{TreeSequenceInitialSegment}, together with the outside vertex
recorded by \noderef{TreeSequenceAdjoinEdge}, identify their complete
first new pair.  Hence the whole nonparallel seed $T_n+f_n$ lies in $A$,
and $e_1$ is not the SE/RE stored-seed exception.

In PE, $T_n$ is an ordered prefix of the hierarchy base core and hence
of $A$; in SE/RE the exact seed $T_n+f_n$ is an ordered prefix of $A$ by
the preceding paragraph.  Also $A\preceq T\preceq T_{n+1}$ by the
hierarchy nesting and (4a).  Expand the certified closure as its finite
\noderef{TashkinovAugmentationSequence}.  Each
\noderef{TashkinovAugmentation} raises the length by exactly one and
retains all earlier vertex and exact edge entries.  The sequence must
therefore contain unique members of lengths $|A|$ and $|T|$; since both
members and $A,T$ are ordered prefixes of $T_{n+1}$ of those respective
lengths, entrywise equality in \noderef{TreeSequenceInitialSegment}
identifies them with $A$ and $T$.  As $|T|=|A|+1$, the intervening
augmentation is precisely the exact addition $(e_1,y_1)$.  Write
\[
                         \varphi_n(e_1)=\gamma.                \tag{4}
\]
The augmentation witness gives
$\gamma\in\overline\varphi_n(A)$.  The terminal `minusClosed' field of
\noderef{GoodHierarchy}, after expanding
\noderef{TreeColorSetMinusClosed}, forces
\[
 \gamma\in
   \bigcup_{\eta\in D_{n,q}}\Gamma^{q-1}_\eta;               \tag{5}
\]
otherwise the boundary edge $e_1$ of the missing color $\gamma$ would
violate that field.

Let $A^-$ be the starred stage preceding $A$, including the core
interpretation when $q=1$, and let $i=q-1\in\operatorname{Fin}(q)$.
The `uniqueRefresh' field of \noderef{GoodHierarchy} supplies
$\eta_g\in D_{n,q}$ such that
\[
 \Gamma^q_{\eta_g}\subseteq
 \bigcup_{x\in V(A)-V(A^-)}\overline\varphi_n(x).             \tag{6}
\]
The `twoColors' field says that this refreshed pair has two members.
Choose $\theta\in\Gamma^q_{\eta_g}$.  Equation (6) supplies an actual
$x\in V(A)-V(A^-)$ missing $\theta$.  By elementarity (2), no vertex of
$A^-$ also misses $\theta$.  Thus $\theta$ is newly missing above the
last divider, and \noderef{GoodHierarchyNewMissingClosed} gives
\[
 \theta\in\overline\varphi_n(A),\qquad
 \partial_{\varphi_n,\theta}(A)=\varnothing.                 \tag{7}
\]
For every $\eta\in D_{n,q}$,
\noderef{HierarchyStoredColorsAntitone} places $\eta$ in the preceding
stored set.  The avoidance field of \noderef{GoodHierarchy} then puts
each color of the union in (5) in the preceding
\noderef{HierarchyStageMissingColors}, hence supplies a missing witness
in $A^-$.  The preceding paragraph says $\theta$ has no such witness,
so
\[
                         \theta\ne\gamma.                    \tag{8}
\]

We first rule out $\theta=\alpha$.  Under that equality, apply
\noderef{ChangeLemma} to $\theta,\gamma$.  Its penultimate prefix is
exactly $A$ because $\ell=1$; (2), (4), and (7) give both missing-color
premises and closedness.  Its conditional ordered-unused premise is
vacuous because both colors are missing on $A$.  Conclusion (iii) makes
$\theta,\gamma$ $A$-interchangeable.  Missing witnesses in $A$ lie in
an isolated full component if they coincide, and otherwise properness,
\noderef{TwoColorIncidenceDegreeAtMostTwo}, and
\noderef{TwoColorComponentSpanningPath} make them the two ends of a
spanning path.  The component containing the exact edge $e_1$ also
meets $A$ and has the additional end $y_1$, where $\theta=\alpha$ is
missing by (3).  The two components are distinct in either case,
contradicting \noderef{TreeInterchangeable}.  Therefore
\[
                         \theta\ne\alpha.                    \tag{9}
\]

Apply \noderef{step1} to the last positive hierarchy stage with first
color $\theta$ and second color $\alpha$.  Its hypotheses are the stated
lower-rung package, fewer-divider induction, large estimate, maximum
property and positivity packaged by $C$, and (7)--(9).  It follows that
$\theta$ and $\alpha$ are $A$-interchangeable.  Choose their missing
witnesses in $A$ as above.  The full
$(\theta,\alpha)$-component $Q$ through $y_1$ cannot meet $A$: equality
with the component determined by the two witnesses would either equate
a nontrivial component with an isolated vertex or add the third end
$y_1$.  Properness and
\noderef{TwoColorIncidenceDegreeAtMostTwo} make $Q$ a possibly singleton
\noderef{TwoColorPathComponent}.

Put
\[
 \sigma_n=\noderef{TwoColorComponentSwap}
       (G,\varphi_n,\theta,\alpha,\operatorname{inl}y_1).     \tag{10}
\]
We record both certified transports needed later.  First apply
\noderef{GoodHierarchyPrefixRepackage} at the exact stage $A$.  It gives
an original same-rung maximum ETT $\mathcal E_A$ on $A$, with the
original certified series and a good retained hierarchy having $q-1$
dividers.  Since the full support of $Q$ is outside $A$, (10) is an
\noderef{OutsideKempeRecoloring} of this prefix.  Applying
\noderef{HierarchyLEMStable} gives a corresponding replacement
$\mathcal E_A^\sigma$, an actual base $\mathcal B_A^\sigma$, an actual
hierarchy $K_A$, and a good witness $g_A$.  It also gives maximum
property, properness, unique-uncoloredness,
\noderef{StableColoring}, \noderef{TreeColoringInvariant}, and
\noderef{SameGoodHierarchy} with the retained original prefix
hierarchy.  The invocation includes the singleton possibility for $Q$
and recolors only its exact edge-nodes.

Second, use positive-divider Case 1 of \noderef{StableChange}, with
colors $\theta,\alpha$ and root $y_1$.  Its full-tail penultimate prefix
is $A$, and (2), (3), (7), (9), the path-component certificate, and the
component disjointness just proved supply its case hypotheses.  The
conclusion gives
\[
 \begin{gathered}
  \sigma_n\text{ weakly stable on }A,
  \qquad \noderef{ModColoring}(T,\sigma_n),                   \tag{11}\\
  A\text{ carrying the transported maximum prefix hierarchy and being}
  \left(\bigcup_{\eta\in D_{n,q}}\Gamma^{q-1}_\eta\right)^-
  \text{-closed under }\sigma_n.                             \tag{12}
 \end{gathered}
\]
It preserves every retained $D$-set and every paper Gamma pair.

We now construct the full same-rung replacement explicitly.  The
edge-node $\operatorname{inr}e_1$ cannot belong to $Q$, since its
adjacency to its endpoint in $A$ would put that vertex-node in $Q$ by
\noderef{TwoColorComponent}.  Its color $\gamma$ differs from $\theta$
by (8), and from $\alpha$ because it is incident with $y_1$, where
$\alpha$ is missing.  The cases in
\noderef{TwoColorComponentSwap} therefore give
\[
                         \sigma_n(e_1)=\gamma.                \tag{13}
\]
The same support argument fixes every edge incident with $A$, so every
missing-color set on $A$ is preserved.

Return to the certified last positive step selected above from
\noderef{TashkinovSeries}; its generating coloring is $\varphi_n$ and
its exact closure output contains $T$ by
\noderef{ExtendedTashkinovTree}.  In the PE case, truncate its certified
augmentation sequence $T_n\longrightarrow T_{n+1}$ at the already
identified member $A$.  In the SE/RE cases, retain the certified exact
seed $T_n+f_n$, which is a prefix of $A$, and truncate the certified
sequence $T_n+f_n\longrightarrow T_{n+1}$ at $A$.  Every augmentation
retained before $A$ uses an exact new edge incident with its old endpoint
in $A$ and a color missing at an old vertex in $A$.  The edge and vertex
clauses of \noderef{TreeColoringInvariant} from (11) preserve,
respectively, that exact edge color and that missing-color witness.
Induction through \noderef{TashkinovAugmentationSequence} therefore
gives the same exact route to $A$ under $\sigma_n$.  Append the exact
augmentation $(e_1,y_1)$, using (13) and the preserved original
$\gamma$-missing witness on $A$.  Thus in PE there is a certified route
$T_n\longrightarrow T$, while in SE/RE there is a certified route
$T_n+f_n\longrightarrow T$ retaining both the same edge identity $f_n$
and $F_n=\{f_n\}$.

These are precisely the two eligibility alternatives in
\noderef{ModColoring}; moreover (4a) supplies their required exact
$T_n\prec T$ ordered prefix.  The universal clause of (11) consequently
supplies a full \noderef{CorrespondingExtendedTashkinovTree} certificate
$\mathcal C_T$ on this same ordered $T$; write its replacement ETT as
$\mathcal E_T^\sigma$.  This is an actual ETT presentation, not merely a
tree inclusion.  Indeed, the `certified' field of its replacement
\noderef{TashkinovSeries} supplies the actual SE/PE/RE certificate at
its last positive index, while the fields of its replacement
\noderef{ExtendedTashkinovTree} give
\[
 T_n^\sigma\prec T\preceq T_{n+1}^\sigma,
 \qquad \sigma_n=(\varphi^{\rm replacement})_n.              \tag{13a}
\]
The series-presentation field of $\mathcal C_T$ identifies that ETT
series with its certified replacement series.  Its rung agreement,
positive-ladder-tree agreement, raw stored-set agreement, exact
connecting-edge-set agreement, and extension-kind agreement then
identify $T_n^\sigma$ with the original ordered $T_n$ and match the
replacement last constructor to the selected original PE, SE, or RE
constructor.  Thus (13a) and the two transported routes above justify
the replacement terminal presentation, including the retained exact
nonparallel seed rather than only its endpoint pair.  By
\noderef{MaximumPropertyTransport}, $\mathcal E_T^\sigma$ has maximum
property.

The prefix repackage used the original series.  Consequently the
corresponding certificates for $\mathcal E_A^\sigma$ and
$\mathcal E_T^\sigma$, after transporting their common rung equality,
identify every positive ladder tree with the corresponding original
ladder tree, every raw stored set and exact connecting-edge set with the
original one, and every extension kind with the original kind.  In
particular, their terminal ladder trees are the same ordered $T_n$ and
their terminal kinds and connecting-edge sets agree.

Construct a hierarchy base $\mathcal B_T^\sigma$ for the full
replacement as follows.  If the common terminal kind is PE, eliminate
the two replacement PE constructors.  Each
\noderef{MaximumDefectiveChoice} has an exact singleton connecting edge
joining its chosen support in $T_n$ to its other endpoint outside
$T_n$.  Equality of the exact connecting-edge sets makes these edge
identities equal.  Equality of their unordered endpoint pairs, together
with the inside/outside orientation, makes the two supports equal; the
same comparison with the original PE constructor identifies this with
the original support.  The parallel fields of the old base and
$\mathcal B_A^\sigma$ locate that support at the same unique position in
the ordered \noderef{TreeSequence} $T_n$.  Copy that position, the
closure, and the exact ordered-union core into
$\mathcal B_T^\sigma$.

If the common terminal kind is not PE, the nonparallel field of
$\mathcal B_A^\sigma$ identifies its core with the common terminal tree
$T_n$; use this as the full replacement core and transport the unused
support position and closure.  In both cases the terminal-missing-color
field is reflexive because the base coloring and generating coloring are
$\sigma_n$.  The PE data just constructed verify the parallel field,
and in the other case the nonparallel field is immediate.  This checks
every field of \noderef{HierarchyBase}.  The
\noderef{SameGoodHierarchy} relation for the prefix also identifies this
core, and the closure in the PE case, with the original base data.

Define $K_T$ by copying from $H$ the divider count $q$, every exact
ordered stage, and every divider position.  Its base-stage equation uses
the common terminal ladder tree; its terminal-stage equation is the
same ordered $T$.  Nestedness, strictness, divider positivity,
divider increase, and exact divider-prefix equations are the
corresponding fields of \noderef{TashkinovHierarchy}.  The
`coreInsideFirst' field transports along the base-core equality just
proved.  Hence $K_T$ is an actual hierarchy of the replacement, with no
dependent field borrowed from the separately chosen $K_A$.

Give $K_T$ the same Gamma function as $C.\mathit{good}$.  We now derive,
rather than assume, equality of its indexed stored sets.  Raw stored-set
agreement in $\mathcal C_T$, equality of the two terminal ladder trees,
and the terminal \noderef{StableColoring} vertex clause make the
computed terminal stored sets equal after expanding
\noderef{TashkinovStoredColors}: the bounded unions of raw slots agree,
and the sets of colors subtracted as missing on $T_n$ agree.  Denote this
common computed set by $D_n$.

Since $q>0$, the weak stability in (11), expanded through
\noderef{HierarchyWeaklyStable}, contains
\noderef{TreeColoringInvariant} on $A$.  The copied stages of $K_T$ and
$H$ are identical ordered trees, and every ordinary or starred stage at
an index $j\leq q$ is a prefix of $A$.  The vertex clause of that
invariant therefore gives
\[
 \overline{\sigma_n}(\noderef{HierarchyStarredStage}(K_T,j))
 =\overline{\varphi_n}(\noderef{HierarchyStarredStage}(H,j)). \tag{13b}
\]
For each transported index $j$, expand the directly cited definition
\noderef{HierarchyStoredColorsAt} on both sides:
\[
 \begin{aligned}
 D^\sigma_{n,j}
   &=D_n^\sigma-
     \overline{\sigma_n}(\noderef{HierarchyStarredStage}(K_T,j)),\\
 D_{n,j}
   &=D_n-
     \overline{\varphi_n}(\noderef{HierarchyStarredStage}(H,j)).
 \end{aligned}
\]
The computed-terminal equality $D_n^\sigma=D_n$ and (13b) make these
two set differences equal.  Hence
\[
 D^\sigma_{n,j}=D_{n,j}\quad(0\leq j\leq q),                 \tag{14}
\]
The same vertex invariant gives equality of every corresponding
\noderef{HierarchyStageMissingColors} set.

The segment-color sets agree as well.  For $j<q$, every displayed edge
lies in $A$ and is fixed by the edge clause of the invariant.  For
$j=q$, the only possible displayed edge above $A$ is $e_1$, fixed by
(13).  The first-missing cutoff is unchanged for $\theta$ and $\alpha$
because their first missing vertices already lie in $A$ and the vertex
clause preserves them; all other colors are untouched by
\noderef{TwoColorComponentSwap}.  Therefore every paper-relevant
\noderef{HierarchySegmentColors} set agrees.                   \tag{15}

Now verify the five \noderef{GoodHierarchy} fields.  `twoColors' and
`disjointPairs' copy after (14) rewrites every quantified stored index.
For `avoidance', fix $j$ and
$\eta\in D^\sigma_{n,j}$; equation (14) makes this precisely the old
membership $\eta\in D_{n,j}$, so the old avoidance field applies, while
the stage-missing equality and (15) identify its two target sets.
For `uniqueRefresh' at $i$, equation (14) at $i+1$ identifies both the
distinguished stored-color membership and every competing stored index;
the copied ordered stage difference, its vertex missing sets, and the
Gamma pairs then transport both existence and uniqueness from the old
field.  In the PE case, equation (14) at $j=0$ identifies the stored set
and therefore the stored-indexed Gamma union in `parallelBase'; together
with the proved equalities of $T_n$, support position, closure, and core,
this is exactly the field of $g_A$.  In the non-PE case its antecedent is
false.

For `minusClosed' at an index before the last divider, equation (14) at
$i+1$ identifies the stored set in its control union, and the copied
Gamma data make the field exactly the corresponding field of $g_A$.
At the last divider, including when $q=1$, equation (14) at $j=q$
identifies the replacement control union with the $D_{n,q}$-indexed
union in (12), so (12) is exactly the required final `minusClosed'
field.  Thus all five fields are verified.  Finally, the definition of
\noderef{SameGoodHierarchy} asks for the already proved same stages and
dividers, equality of every transported indexed stored set, and equality
of Gamma values on that common set.  Equation (14) supplies its stored-set
clause, and the copied Gamma function supplies its remaining clause.
This constructs the good witness $g_T$ and proves the claimed
\noderef{SameGoodHierarchy} relation with the original hierarchy.

The swap makes $\theta$ missing at $y_1$, while it preserves the missing
$\theta$ at the vertex $x\in A$ from (6).  Since $y_1\notin V(A)$, the
full replacement tree is not elementary.  The rung agreement in
$\mathcal C_T$, maximum property, the constructed base and good
hierarchy, and positivity therefore form an actual
\noderef{HierarchyCounterexample} $C_\sigma$ on the same ordered $T$.

Its measure equals that of $C$: the ordered tree, divider count, full
stage family, and therefore both coordinates of
\noderef{HierarchyCounterexampleMeasure} are unchanged.  To transfer
minimality, take any $q$-divider counterexample with the same rung as
$C_\sigma$ and compose its rung equality with that of $\mathcal C_T$;
it is eligible for the original minimality hypothesis.  A smaller
measure than that of $C_\sigma$ would, after the displayed measure
equality, be smaller than that of $C$, contradicting the hypothesis.
The tail still has length one and its exact penultimate prefix is $A$;
vertex missing-set invariance transfers (2).  The common terminal ladder
tree and terminal stability preserve $|\overline\varphi_n(T_n)|$, so
the large estimate transfers.  The lower-rung package is unchanged and
the fewer-divider hypothesis rewrites along the divider-count equality.
Properness and the exact one-uncolored-edge condition are supplied by
(11).  Thus every ambient hypothesis of \noderef{ChangeLemma} holds for
$C_\sigma$.

Under $\sigma_n$, both $\theta$ and $\gamma$ are missing on $A$.
Equations (5), (7), (12), and the pointwise invariance show that
$\theta$ is outside the reserved union and closed on $A$, after
expanding \noderef{TreeColorSetMinusClosed} and
\noderef{TreeColorSetClosed}.  Apply conclusion (iii) of
\noderef{ChangeLemma} for $C_\sigma$ to $\theta,\gamma$; its conditional
ordered-unused premise is vacuous because both colors are missing on
$A$.  It makes them $A$-interchangeable.

Finally, $y_1$ misses $\theta$ after the swap, while the unchanged exact
edge $e_1$ has color $\gamma$, so its full
$(\theta,\gamma)$-component meets $A$.  Missing witnesses for the two
colors in $A$ determine an isolated component if they coincide, or a
spanning path with exactly those two ends if they differ, by
\noderef{TwoColorComponentSpanningPath}.  It cannot be the component
with additional end $y_1$, by
\noderef{TwoColorIncidenceDegreeAtMostTwo}.  Hence at least two path
components, including isolated components, meet $A$, contradicting
\noderef{TreeInterchangeable}.  This is the required contradiction.
\end{proof}
